% =============================================================================
% capitolo1_contesto.tex
% Capitolo 1: Contesto scientifico e tecnologico
%
% LINEE GUIDA VITALI:
%  Questo capitolo risponde alla domanda: "Perche' il problema non e' gia' risolto?"
%
%  A) STATO DELL'ARTE:
%     Descrivi gli approcci esistenti. Per ognuno: cosa fa bene, cosa fa male,
%     perche' non e' sufficiente.
%     CITAZIONI OBBLIGATORIE: ogni affermazione su lavori altrui vuole \cite{XYZ99}.
%     Esempio: "Come dimostrato da \citet{Tiz20}, l'approccio X ha il limite di..."
%
%  B) NUOVI STRUMENTI / ALGORITMI ABILITANTI:
%     Se il tuo lavoro sfrutta tecnologie recenti (LLM, nuovi framework, ecc.),
%     spiegale qui con sufficiente dettaglio da giustificare le scelte del Cap. 2.
%
%  C) IL GAP:
%     Concludi con una sintesi di cosa manca nello stato dell'arte
%     e come il tuo lavoro si posiziona per colmarlo.
%
%  REGOLA D'ORO: Non fare affermazioni senza fonte. Cita sempre con \cite{XYZ99}.
%  LUNGHEZZA SUGGERITA: 15-25 pagine.
% =============================================================================

% --- A. STATO DELL'ARTE ---
% Descrivi i principali approcci esistenti al tuo problema.
% Per ogni approccio: punti di forza, limiti, perche' non basta.

% --- B. TECNOLOGIE E STRUMENTI ABILITANTI ---
% Descrivi le tecnologie chiave su cui si basa il tuo prototipo.
% Spiega perche' le hai scelte rispetto ad alternative.

% --- C. GAP E MOTIVAZIONE ---
% Sintetizza: cosa manca nello stato dell'arte?
% Questa sezione prepara il lettore al Capitolo 2.

Scrivi qui il contenuto del capitolo 1 sul contesto scientifico.
